\section{Background}

A \emph{compressed static function} (CSF)~\cite{hreinsson2009storing} is a data structure
that stores a mapping $f : X \to \Sigma$ from a finite key set $X$ of size $n$ to an
alphabet $\Sigma$, using space per key close to the empirical entropy of the output
values. The core construction encodes each output value with a variable-length
prefix-free codeword and builds a bit array $A$ such that XOR-ing $k$ positions of
$A$, at positions determined by hash functions of the key, recovers the codeword.
This reduction from a function-storage problem to a system of linear equations over
$\mathbb{Z}/2\mathbb{Z}$ (specifically a $k$-XORSAT instance) yields a bit array of
size approximately
\[
    \delta_k \sum_{x \in X} |c(f(x))| \text{bits},
\]
where $c(\cdot)$ is the chosen prefix-free code. 

When the code matches the empirical distribution of outputs, the average codeword
length equals the empirical entropy $H_0$, so the bit array consumes $\approx \delta_k
n H_0 \text{bits}$ in total. Beyond the bit array, one must also store the \emph{codebook}:
the Huffman decoding table and the mapping from codewords to elements of $\Sigma$.

